\documentclass[11pt,a4paper]{article}
\usepackage[utf8]{inputenc}
\usepackage[T1]{fontenc}
\usepackage{geometry}
\usepackage{amsmath,amssymb}
\usepackage{booktabs}
\usepackage{longtable}
\usepackage{enumitem}
\usepackage{xcolor}
\usepackage{tcolorbox}
\usepackage{listings}
\usepackage{hyperref}
\usepackage{fancyhdr}
\usepackage{titlesec}
\usepackage{array}
\usepackage{multirow}
\usepackage{float}
\usepackage{mdframed}
\usepackage{colortbl}
\usepackage{tabularx}

\geometry{margin=1in}
\hypersetup{colorlinks=true,linkcolor=blue,urlcolor=blue}

\definecolor{critical}{RGB}{220,53,69}
\definecolor{warning}{RGB}{255,193,7}
\definecolor{success}{RGB}{40,167,69}
\definecolor{info}{RGB}{23,162,184}
\definecolor{codeblue}{RGB}{0,102,204}
\definecolor{codegray}{RGB}{128,128,128}
\definecolor{codebg}{RGB}{248,248,248}
\definecolor{nunchiblue}{HTML}{1a73e8}
\definecolor{nunchidark}{HTML}{202124}
\definecolor{cooperative}{RGB}{106,27,154}
\definecolor{keeper}{RGB}{2,119,189}
\definecolor{operator}{RGB}{230,81,0}
\definecolor{managed}{RGB}{46,125,50}

\newtcolorbox{criticalbox}{colback=critical!10,colframe=critical,title=\textbf{CRITICAL}}
\newtcolorbox{warningbox}{colback=warning!10,colframe=warning!80!black,title=\textbf{WARNING}}
\newtcolorbox{infobox}{colback=info!10,colframe=info,title=\textbf{INFO}}
\newtcolorbox{successbox}{colback=success!10,colframe=success,title=\textbf{SUCCESS CRITERIA}}
\newtcolorbox{contractbox}{colback=codebg,colframe=codeblue,title=\textbf{Contract Interface}}
\newtcolorbox{flowbox}{colback=gray!5,colframe=gray!50!black,title=\textbf{Process Flow}}
\newtcolorbox{cooperativebox}{colback=cooperative!8,colframe=cooperative,title=\textbf{House Cooperative}}
\newtcolorbox{keeperbox}{colback=keeper!8,colframe=keeper,title=\textbf{Permissionless Keeper}}
\newtcolorbox{operatorbox}{colback=operator!8,colframe=operator,title=\textbf{Privileged Operator}}
\newtcolorbox{managedbox}{colback=managed!8,colframe=managed,title=\textbf{Managed Infrastructure}}
\newtcolorbox{designbox}[1][]{colback=nunchiblue!5,colframe=nunchiblue,title=\textbf{Design Decision: #1}}

% CLI-specific box
\newtcolorbox{clibox}[1][]{colback=nunchidark!5,colframe=nunchidark,title=\textbf{CLI: #1}}

\lstset{
    basicstyle=\ttfamily\small,
    backgroundcolor=\color{codebg},
    keywordstyle=\color{codeblue}\bfseries,
    commentstyle=\color{codegray},
    stringstyle=\color{success!80!black},
    breaklines=true,
    frame=single,
    rulecolor=\color{codegray},
    tabsize=2
}

\pagestyle{fancy}
\fancyhf{}
\rhead{Nunchi Trade -- Agent CLI Jobs}
\lhead{Specification v1.0}
\rfoot{Page \thepage}

\title{\textbf{Agent CLI: Job Integration \& Assignment Routing}\\[0.5em]
\large Unified CLI for Perpetual Agent Jobs\\[0.3em]
\small Companion to: Perpetual Agent Jobs Specification v1.0}
\author{Nunchi Trade}
\date{March 2026}

\begin{document}

\maketitle

\begin{center}
\fbox{\parbox{0.9\textwidth}{
\textbf{Scope.} This document specifies how the \texttt{hl} CLI integrates with the
eight perpetual agent jobs defined in the Perpetual Agent Jobs Specification v1.0.
It covers the \texttt{hl jobs} command group, the engine hierarchy that dispatches
each job to the correct execution runtime, strategy interfaces for each job category,
event subscription, custody enforcement at the CLI layer, and the phased migration
plan from current-state scripts to a unified CLI-driven job system.
}}
\end{center}

\tableofcontents
\newpage

%==============================================================================
\section{Introduction and Motivation}
%==============================================================================

The Nunchi agent runtime currently spans two disconnected execution surfaces:

\begin{enumerate}[leftmargin=2em]
  \item \textbf{The \texttt{hl} CLI} handles free-form Hyperliquid trading strategies.
        A user runs \texttt{hl run avellaneda\_mm} to start a market-making loop that
        polls the relay, generates quotes via a pluggable \texttt{BaseStrategy}, and
        submits sealed decisions through the TEE-work commit--reveal pipeline.
  \item \textbf{Standalone scripts} (\texttt{scripts/run\_agent.py},
        \texttt{scripts/run\_relay.py}) bootstrap the cooperative market-making
        infrastructure separately. Keeper services such as oracle updates, funding
        settlement, and liquidation flagging are handled by ad-hoc cron jobs or
        manual invocations with no shared lifecycle management.
\end{enumerate}

The Perpetual Agent Jobs Specification formalizes eight on-chain roles across four
categories---\textcolor{keeper}{\textbf{Keeper}}, \textcolor{operator}{\textbf{Operator}},
\textcolor{cooperative}{\textbf{Cooperative}}, and \textcolor{managed}{\textbf{Managed}}---each
with explicit trigger conditions, custody constraints, staking requirements, and
reward mechanisms. This document specifies how the CLI becomes the \textbf{unified
interface} for all eight job types, from permissionless keepers to TEE-attested
cooperative agents.

Three requirements drive the design:

\begin{enumerate}[leftmargin=2em]
  \item \textbf{Job integration.} Register, run, monitor, and claim rewards for any
        of the 8 job types through a single CLI command group (\texttt{hl jobs}).
        The CLI manages the full on-chain lifecycle: staking, heartbeats, reward
        claims, and deregistration.
  \item \textbf{Assignment routing.} Automatically dispatch each job to the correct
        execution engine based on its category. A keeper job routes to the
        \texttt{KeeperEngine} (event-driven, stateless transaction submission);
        a cooperative job routes to the \texttt{CooperativeEngine} (relay-polling,
        commit--reveal rounds); a managed job routes to the \texttt{ManagedEngine}
        (periodic evaluation with capital actions). The user does not need to know
        which engine handles their job---the factory resolves it from the job
        definition.
  \item \textbf{Extensibility.} Support new job types without modifying the routing
        or engine infrastructure. Adding a new keeper (e.g., a cross-chain bridge
        relayer) requires only a new \texttt{KeeperStrategy} implementation and a
        registry entry. The engine, event subscription, and custody layers remain
        unchanged.
\end{enumerate}

\begin{infobox}
This spec is a companion to the Perpetual Agent Jobs Specification v1.0. That document
defines the jobs themselves (triggers, custody, rewards, slashing). This document
defines the CLI surface and execution infrastructure that runs them.
\end{infobox}

%==============================================================================
\section{CLI Command Design}
%==============================================================================

All job-related commands live under the \texttt{hl jobs} command group. Each command
maps to a specific phase of the job lifecycle.

\begin{clibox}[hl jobs list]
\begin{lstlisting}
$ hl jobs list

 ID                 NAME              CATEGORY     TRIGGER          MIN STAKE
 oracle_updater     Oracle Updater    KEEPER       NewBlock         10 HYPE
 funding_keeper     Funding Keeper    KEEPER       NewBlock         10 HYPE
 liq_flagger        Liq. Flagger      KEEPER       OracleUpdate     10 HYPE
 liq_executor       Liq. Executor     OPERATOR     PositionFlagged  50 HYPE
 tpsl_agent         TP/SL Agent       OPERATOR     OracleUpdate     25 HYPE
 market_maker       Market Maker      COOPERATIVE  RoundStart       100 HYPE
 abm_agent          ABM Agent         COOPERATIVE  OracleUpdate     100 HYPE
 glv_manager        GLV Manager       MANAGED      Periodic+Event   50 HYPE
\end{lstlisting}
Lists all job types from the local registry, showing category, trigger type, and minimum stake.
\end{clibox}

\begin{clibox}[hl jobs info]
\begin{lstlisting}
$ hl jobs info oracle_updater

  Job:       Oracle Updater (oracle_updater)
  Category:  KEEPER (Permissionless)
  Trigger:   NewBlock
  Custody:   OracleManager.updatePriceFeeds only
  Min Stake: 10 HYPE
  Reward:    0.002 HYPE per successful update
  Slashing:  Liveness (miss 100 consecutive blocks)
  TEE:       Not required
  Status:    Not registered
\end{lstlisting}
Shows full details for a specific job: trigger conditions, custody constraints,
staking and reward parameters, slashing conditions, and current registration status.
\end{clibox}

\begin{clibox}[hl jobs register]
\begin{lstlisting}
$ hl jobs register oracle_updater --stake 15

  Registering for job: Oracle Updater
  Staking: 15 HYPE (min: 10 HYPE)
  TEE attestation: not required
  Transaction: 0xabc123...
  Status: REGISTERED
\end{lstlisting}
Registers the caller for a job on-chain. Stakes the specified amount (must meet
minimum). If \texttt{-{}-tee} is passed, generates and submits a TEE attestation
as part of registration.
\end{clibox}

\begin{clibox}[hl jobs run]
\begin{lstlisting}
$ hl jobs run oracle_updater --config jobs/oracle.yaml --mainnet

  Loading config: jobs/oracle.yaml
  Job: oracle_updater -> KeeperEngine
  Subscribing to: NewBlock events
  Custody policy: OracleManager.updatePriceFeeds
  Engine started. Press Ctrl+C to stop.
  [12:00:01] Block 1234: staleness 47s > 45s, submitting update...
  [12:00:01] tx 0xdef456 submitted (gas: 142000)
  [12:00:13] Block 1235: staleness 1s, skipping.
\end{lstlisting}
Starts the job engine. The CLI loads the YAML config, looks up the job definition
in the registry, validates requirements (role, TEE, minimum stake), creates the
appropriate engine via \texttt{JobEngineFactory}, and enters the execution loop.
Supports \texttt{-{}-dry-run} for local simulation without submitting transactions.
\end{clibox}

\begin{clibox}[hl jobs status]
\begin{lstlisting}
$ hl jobs status

  RUNNING JOBS:
  ID                ENGINE          UPTIME    HEARTBEAT   REWARDS
  oracle_updater    KeeperEngine    2h 14m    12s ago     1.24 HYPE
  market_maker      CoopEngine      1h 02m    3s ago      8.71 HYPE
\end{lstlisting}
Shows running jobs, their engine type, uptime, last heartbeat, and accumulated rewards.
Use \texttt{-{}-job-id <id>} to filter to a specific job.
\end{clibox}

\begin{clibox}[hl jobs claim]
\begin{lstlisting}
$ hl jobs claim oracle_updater

  Claiming rewards for: Oracle Updater
  Accumulated: 1.24 HYPE
  Transaction: 0x789abc...
  Claimed: 1.24 HYPE -> 0xYourAddress
\end{lstlisting}
Claims accumulated rewards for a job from the on-chain JobRegistry contract.
\end{clibox}

\begin{clibox}[hl jobs deregister]
\begin{lstlisting}
$ hl jobs deregister oracle_updater

  Deregistering from: Oracle Updater
  Unstaking: 15 HYPE (cooldown: 7 days)
  Transaction: 0xfed321...
  Status: DEREGISTERED (stake unlocks at block 98765)
\end{lstlisting}
Deregisters from a job and initiates the unstaking cooldown. The staked amount
is returned after the cooldown period.
\end{clibox}

\begin{designbox}[Why \texttt{hl jobs} is separate from \texttt{hl run}]
The existing \texttt{hl run <strategy>} command starts a free-form trading strategy
with no on-chain lifecycle. There is no registration, no staking, no heartbeat
obligation, no slashing risk, and no custody enforcement. The user simply picks a
strategy and runs it against Hyperliquid.

Jobs are fundamentally different: they have an \textbf{on-chain lifecycle}
(registration, staking, heartbeats, slashing, custody enforcement) that trading
strategies do not need. Merging these into a single command would force trading
users to deal with staking and registration, and force job operators to skip
lifecycle steps they actually need. The \texttt{hl jobs} command group makes
the lifecycle explicit while reusing the same strategy infrastructure underneath.

A cooperative job (\texttt{hl jobs run market\_maker}) and a free-form strategy
(\texttt{hl run avellaneda\_mm}) may even use the same \texttt{BaseStrategy}
subclass, but the job version wraps it with heartbeats, custody validation, and
reward tracking.
\end{designbox}

\subsection{Job Configuration (YAML)}

Each job is configured via a YAML file passed to \texttt{hl jobs run -{}-config}.

\begin{keeperbox}
\textbf{Example: Keeper job configuration (oracle\_updater)}
\begin{lstlisting}[language={}]
job_id: oracle_updater
agent_id: agent-oracle-1
chain_rpc: https://rpc.nunchi.trade
event_bus_ws: wss://events.nunchi.trade/ws
strategy: jobs.keepers.oracle:OracleKeeperStrategy
strategy_params:
  staleness_threshold_s: 45
  deviation_threshold_pct: 0.1
  pyth_hermes_url: https://hermes.pyth.network
mainnet: true
data_dir: data/jobs/oracle
\end{lstlisting}
\end{keeperbox}

\begin{cooperativebox}
\textbf{Example: Cooperative job configuration (market\_maker)}
\begin{lstlisting}[language={}]
job_id: market_maker
agent_id: agent-mm-1
relay_url: http://relay.nunchi.trade:8080
strategy: strategies.avellaneda_mm:AvellanedaStoikovMM
strategy_params:
  gamma: 0.1
  k: 4.0
  base_size: 0.05
tee_enabled: true
stake_amount: 100.0
risk:
  max_position_qty: 0.5
  max_leverage: 2.0
  max_daily_drawdown_pct: 2.0
data_dir: data/jobs/mm
\end{lstlisting}
\end{cooperativebox}

%==============================================================================
\section{Architecture --- Engine Hierarchy}
%==============================================================================

The engine hierarchy is the core dispatch mechanism. Each job category maps to
exactly one engine type. All engines implement the same \texttt{JobEngine} abstract
base class, ensuring uniform lifecycle management regardless of the underlying
execution model.

\begin{contractbox}
\textbf{Engine Hierarchy}
\begin{lstlisting}[language={}]
JobEngine (ABC)
  start(job_def, config) -> None
  stop() -> None
  heartbeat() -> None
  status() -> JobStatus
  |
  +-- KeeperEngine          # KEEPER + OPERATOR jobs
  |     EventSubscriber -> condition check -> CustodyGuard
  |       -> tx submit -> heartbeat
  |
  +-- CooperativeEngine     # COOPERATIVE jobs (wraps AgentClient)
  |     Relay poll -> strategy tick -> seal -> commit
  |       -> reveal -> feedback -> heartbeat
  |
  +-- ManagedEngine          # MANAGED jobs
        Timer + EventSubscriber -> condition check
          -> CustodyGuard -> tx submit
\end{lstlisting}
\end{contractbox}

\textbf{KeeperEngine} handles both \textcolor{keeper}{Keeper} and
\textcolor{operator}{Operator} jobs. These are event-driven: the engine subscribes
to chain events (new blocks, oracle updates, position flags), evaluates a
\texttt{KeeperStrategy} on each event, and submits transactions when the strategy
returns a non-null result. The \texttt{CustodyGuard} validates every transaction
against the job's custody policy before signing.

\textbf{CooperativeEngine} handles \textcolor{cooperative}{Cooperative} jobs. It
wraps the existing \texttt{AgentClient} from the tee-work-llm repository, adding
heartbeat reporting and reward tracking. The execution model is relay-driven:
poll for a clearing round, run the strategy tick, seal the decision, commit, reveal,
and receive feedback.

\textbf{ManagedEngine} handles \textcolor{managed}{Managed} jobs. It combines
periodic timers with event subscriptions. The \texttt{ManagedStrategy} evaluates
vault state on each tick and returns a list of capital actions (rebalance, harvest,
process withdrawals). Like \texttt{KeeperEngine}, it enforces custody constraints
before transaction submission.

\subsection{Factory Dispatch}

The \texttt{JobEngineFactory} maps job categories to engine types. This is the
single point of routing logic---adding a new category requires only a new match arm.

\begin{contractbox}
\textbf{JobEngineFactory}
\begin{lstlisting}[language=Python]
class JobEngineFactory:
    @staticmethod
    def create(job_def: JobDefinition) -> JobEngine:
        match job_def.category:
            case JobCategory.KEEPER | JobCategory.OPERATOR:
                return KeeperEngine(job_def)
            case JobCategory.COOPERATIVE:
                return CooperativeEngine(job_def)
            case JobCategory.MANAGED:
                return ManagedEngine(job_def)
            case _:
                raise ValueError(
                    f"Unknown job category: {job_def.category}"
                )
\end{lstlisting}
\end{contractbox}

\subsection{Full Dispatch Flow}

\begin{flowbox}
\textbf{\texttt{hl jobs run <job\_id> -{}-config <file>}}
\begin{enumerate}[leftmargin=2em]
  \item Load \texttt{JobConfig} from the YAML configuration file.
  \item Look up \texttt{job\_id} in \texttt{JOB\_REGISTRY} to obtain the
        \texttt{JobDefinition} (category, trigger, custody policy, staking
        requirements).
  \item \textbf{Validate}: check that the caller meets role requirements (e.g.,
        \texttt{OPERATOR\_ROLE} for privileged jobs), TEE requirements (attestation
        present if \texttt{requires\_tee} is true), and minimum stake (on-chain
        registration with sufficient stake).
  \item \texttt{JobEngineFactory.create(job\_def)} returns the appropriate engine
        instance.
  \item \texttt{engine.start(job\_def, config)} initializes event subscriptions,
        loads the strategy, and enters the execution loop.
  \item Register signal handlers (\texttt{SIGINT}, \texttt{SIGTERM}) for graceful
        shutdown.
  \item On signal: \texttt{engine.stop()} flushes pending state, sends a final
        heartbeat, and exits cleanly.
\end{enumerate}
\end{flowbox}

%==============================================================================
\section{Strategy Interfaces}
%==============================================================================

Each engine type consumes a different strategy abstract base class. The strategy
is the user-authored component: it encapsulates the decision logic while the engine
handles lifecycle, event routing, custody, and heartbeats.

\subsection{KeeperStrategy}

\begin{contractbox}
\textbf{KeeperStrategy --- for Keeper and Operator jobs}
\begin{lstlisting}[language=Python]
class KeeperStrategy(ABC):
    """Event-driven, stateless decisions.

    Receives a chain event and context, returns either a
    Transaction to submit or None to skip this event.
    """

    @abstractmethod
    def should_execute(
        self,
        event: ChainEvent,
        context: KeeperContext,
    ) -> Optional[Transaction]:
        """Evaluate the event against current on-chain state.

        Args:
            event: The chain event that triggered evaluation
                   (new block, oracle update, position flag).
            context: Current on-chain state relevant to this
                     job (price ages, margin ratios, etc.).

        Returns:
            A Transaction to submit, or None to skip.
        """
        ...
\end{lstlisting}
\end{contractbox}

\textbf{KeeperContext} contains:
\begin{itemize}[leftmargin=2em]
  \item \texttt{event}: the \texttt{ChainEvent} that triggered evaluation (block number, event type, raw log data).
  \item \texttt{on\_chain\_state}: current state relevant to the job---e.g., oracle price age and deviation for the oracle updater, flagged position list for the liquidation executor, margin ratios for the liquidation flagger.
  \item \texttt{gas\_estimate}: estimated gas cost for the transaction, used by the strategy to decide whether execution is profitable given current reward rates.
\end{itemize}

\subsection{CooperativeStrategy}

\begin{contractbox}
\textbf{CooperativeStrategy --- for House Cooperative jobs}
\begin{lstlisting}[language=Python]
class CooperativeStrategy(BaseStrategy):
    """Extends BaseStrategy with round feedback.

    Inherits the core tick() -> List[StrategyDecision]
    interface from BaseStrategy. Adds on_round_result()
    for post-clearing adaptive learning.
    """

    def on_round_result(
        self,
        result: RoundResult,
    ) -> None:
        """Receive post-clearing feedback.

        Called after the clearing engine resolves the round.
        Contains fill information, PnL attribution, and
        counterparty flow data for adaptive learning.

        Args:
            result: Clearing outcome including fills,
                    realized PnL, and inventory delta.
        """
        pass
\end{lstlisting}
\end{contractbox}

\textbf{CooperativeStrategy} extends the existing \texttt{BaseStrategy} from the
agent runtime. The inherited \texttt{tick()} method receives a \texttt{StrategyContext}
(current position, PnL, market snapshot) and returns a list of \texttt{StrategyDecision}
objects (quotes, hedges). The added \texttt{on\_round\_result()} callback enables
strategies to learn from clearing outcomes---adjusting spread, inventory targets,
or aggressiveness based on realized fills and adverse selection.

\subsection{ManagedStrategy}

\begin{contractbox}
\textbf{ManagedStrategy --- for Managed Infrastructure jobs}
\begin{lstlisting}[language=Python]
class ManagedStrategy(ABC):
    """Periodic evaluation with capital management actions.

    Evaluates vault state on each tick and returns a list
    of capital actions to execute (rebalance, harvest fees,
    process pending withdrawals).
    """

    @abstractmethod
    def evaluate(
        self,
        context: ManagedContext,
    ) -> List[CapitalAction]:
        """Evaluate current state and return actions.

        Args:
            context: Current vault state including balances,
                     pending withdrawals, total assets, and
                     last harvest timestamp.

        Returns:
            List of CapitalAction objects to execute.
            Empty list means no action needed.
        """
        ...
\end{lstlisting}
\end{contractbox}

\textbf{ManagedContext} contains:
\begin{itemize}[leftmargin=2em]
  \item \texttt{vault\_balance}: current idle balance in the vault contract.
  \item \texttt{pending\_withdrawals}: list of withdrawal requests awaiting processing, each with amount and requester.
  \item \texttt{total\_assets}: total assets under management (idle + deployed in strategy).
  \item \texttt{last\_harvest\_timestamp}: Unix timestamp of the last performance fee harvest, used to determine if a new harvest is due.
\end{itemize}

\subsection{Job-to-Strategy Mapping}

\begin{table}[H]
\centering
\small
\begin{tabularx}{\textwidth}{l l l X}
\toprule
\textbf{Job} & \textbf{Strategy ABC} & \textbf{Context Model} & \textbf{Decision Output} \\
\midrule
\rowcolor{keeper!8}
Oracle Updater & KeeperStrategy & KeeperContext (price age, Pyth data) & Transaction (\texttt{updatePriceFeeds}) \\
\rowcolor{keeper!8}
Funding Keeper & KeeperStrategy & KeeperContext (window state) & Transaction (\texttt{settleFundingWindow}) \\
\rowcolor{keeper!8}
Liq.\ Flagger & KeeperStrategy & KeeperContext (positions, margins) & Transaction (\texttt{flagPosition}) \\
\rowcolor{operator!8}
Liq.\ Executor & KeeperStrategy & KeeperContext (flagged positions) & Transaction (\texttt{liquidatePosition}) \\
\rowcolor{operator!8}
TP/SL Agent & KeeperStrategy & KeeperContext (conditional orders, prices) & Transaction (\texttt{placeOrder} REDUCE\_ONLY) \\
\rowcolor{cooperative!8}
Market Maker & CooperativeStrategy & StrategyContext (position, PnL) & List[StrategyDecision] \\
\rowcolor{cooperative!8}
ABM Agent & CooperativeStrategy & StrategyContext + bin context & List[StrategyDecision] \\
\rowcolor{managed!8}
GLV Manager & ManagedStrategy & ManagedContext (vault, withdrawals) & List[CapitalAction] \\
\bottomrule
\end{tabularx}
\caption{Mapping of each job to its strategy interface, context model, and output type.}
\end{table}

%==============================================================================
\section{Event Subscription Model}
%==============================================================================

Engines receive triggers through an \texttt{EventSubscriber} abstraction. This
decouples the engine from the specific event transport---whether that is a WebSocket
connection to the canonical chain event bus or a fallback polling mechanism against
EVM \texttt{eth\_getLogs}.

\begin{contractbox}
\textbf{EventSubscriber ABC}
\begin{lstlisting}[language=Python]
class EventSubscriber(ABC):
    """Abstract event source for job engines."""

    @abstractmethod
    async def subscribe(
        self,
        event_types: List[str],
        callback: Callable[[ChainEvent], Awaitable[None]],
    ) -> None:
        """Subscribe to one or more event types.

        Args:
            event_types: List of event type strings
                         (e.g., ["NewBlock", "OracleUpdate"]).
            callback: Async function called on each event.
        """
        ...

    @abstractmethod
    async def unsubscribe(self) -> None:
        """Cleanly disconnect and release resources."""
        ...
\end{lstlisting}
\end{contractbox}

\subsection{Implementations}

\begin{contractbox}
\textbf{ChainEventSubscriber}
\begin{lstlisting}[language=Python]
class ChainEventSubscriber(EventSubscriber):
    """WebSocket connection to canonical chain event bus.

    Connects to the chain event bus via WebSocket, subscribes
    to requested event types, and dispatches ChainEvent objects
    to the callback. Handles reconnection with exponential
    backoff. Target implementation for L1 mainnet.
    """

    def __init__(self, ws_url: str): ...

    async def subscribe(self, event_types, callback): ...
    async def unsubscribe(self): ...
\end{lstlisting}
\end{contractbox}

\begin{contractbox}
\textbf{LogPollingSubscriber}
\begin{lstlisting}[language=Python]
class LogPollingSubscriber(EventSubscriber):
    """Fallback: polls eth_getLogs for contract events.

    For V1 contracts on HyperEVM before chain event bus is
    available. Polls at a configurable interval (default: 2s),
    converts raw logs into ChainEvent objects, and dispatches
    to the callback. Less efficient but works with any
    EVM-compatible RPC endpoint.
    """

    def __init__(self, rpc_url: str, poll_interval_s: float = 2.0): ...

    async def subscribe(self, event_types, callback): ...
    async def unsubscribe(self): ...
\end{lstlisting}
\end{contractbox}

\subsection{Event Type Mapping}

Each job subscribes to specific event types based on its trigger condition. The
engine resolves the correct event types from the \texttt{JobDefinition} and passes
them to the subscriber.

\begin{table}[H]
\centering
\small
\begin{tabularx}{\textwidth}{l X}
\toprule
\textbf{Job} & \textbf{Event Subscription} \\
\midrule
\rowcolor{keeper!8}
Oracle Updater & \texttt{NewBlock} --- evaluate staleness on every new block \\
\rowcolor{keeper!8}
Funding Keeper & \texttt{NewBlock} + window-end check --- evaluate on each block whether the current funding window has ended \\
\rowcolor{keeper!8}
Liq.\ Flagger & \texttt{OracleUpdate} --- re-evaluate margin ratios after each oracle price update \\
\rowcolor{operator!8}
Liq.\ Executor & \texttt{PositionFlagged} --- act immediately when a position is flagged for liquidation \\
\rowcolor{operator!8}
TP/SL Agent & \texttt{OracleUpdate} --- check conditional order trigger prices after each oracle update \\
\rowcolor{cooperative!8}
Market Maker & \texttt{ClearingRoundStart} (from relay, not chain) --- poll relay for new clearing rounds \\
\rowcolor{cooperative!8}
ABM Agent & \texttt{OracleUpdate} + deviation threshold --- re-evaluate bin placement when price moves beyond threshold \\
\rowcolor{managed!8}
GLV Manager & \texttt{WithdrawRequested} + periodic timer --- respond to withdrawal requests and periodically harvest fees \\
\bottomrule
\end{tabularx}
\caption{Event subscriptions per job. Cooperative jobs receive triggers from the relay rather than the chain event bus.}
\end{table}

\begin{designbox}[Two-Phase Event Transport]
The \texttt{EventSubscriber} ABC abstracts the event transport so that engines are
agnostic to the source. This enables a two-phase rollout:

\textbf{Phase 1 (V1 contracts on HyperEVM):} The \texttt{LogPollingSubscriber}
polls \texttt{eth\_getLogs} against deployed V1 contracts. This is functional but
inefficient---it introduces latency proportional to the polling interval and
generates unnecessary RPC load. It serves as the fallback for early deployments
before the chain event bus is operational.

\textbf{Phase 2 (L1 mainnet):} The \texttt{ChainEventSubscriber} connects to
the canonical event bus via WebSocket. Events are pushed in real-time with sub-second
latency. The engine code is unchanged---only the subscriber implementation swapped
via configuration.

This abstraction also supports testing: a \texttt{MockEventSubscriber} can inject
synthetic events for deterministic engine testing without a live chain.
\end{designbox}

%==============================================================================
\section{Custody Enforcement}
%==============================================================================

Custody enforcement at the CLI layer is defense-in-depth. The on-chain contracts
enforce role checks and selector validation independently, but the CLI-side
\texttt{CustodyGuard} catches mistakes \emph{before} they reach the chain---preventing
wasted gas on transactions that would revert.

\begin{contractbox}
\textbf{CustodyGuard}
\begin{lstlisting}[language=Python]
class CustodyGuard:
    """Validates transactions against a job's custody policy
    before signing and submitting."""

    def __init__(self, policy: CustodyPolicy):
        self._policy = policy
        self._tx_count_per_block: Dict[int, int] = {}

    def validate(self, tx: Transaction) -> bool:
        """Validate a transaction against the custody policy.

        Checks:
          1. tx.to in policy.destinations
             (only allowed contract addresses)
          2. tx.data[:4] in policy.selectors
             (only allowed function selectors)
          3. tx.value <= policy.value_cap
             (native value transfer cap)
          4. Rate limit: txs per block does not exceed
             policy.max_txs_per_block

        Returns True if all checks pass, False otherwise.
        Logs the specific violation on failure.
        """
        # 1. Destination whitelist
        if tx.to not in self._policy.destinations:
            logger.warning(
                f"CustodyGuard: tx.to {tx.to} not in "
                f"allowed destinations"
            )
            return False

        # 2. Function selector whitelist
        selector = tx.data[:4]
        if selector not in self._policy.selectors:
            logger.warning(
                f"CustodyGuard: selector {selector.hex()} "
                f"not in allowed selectors"
            )
            return False

        # 3. Value cap
        if tx.value > self._policy.value_cap:
            logger.warning(
                f"CustodyGuard: tx.value {tx.value} "
                f"exceeds cap {self._policy.value_cap}"
            )
            return False

        # 4. Rate limit
        block = tx.block_number
        count = self._tx_count_per_block.get(block, 0)
        if count >= self._policy.max_txs_per_block:
            logger.warning(
                f"CustodyGuard: rate limit exceeded "
                f"for block {block}"
            )
            return False
        self._tx_count_per_block[block] = count + 1

        return True
\end{lstlisting}
\end{contractbox}

\begin{criticalbox}
Custody enforcement is defense-in-depth. The on-chain contracts also enforce role
checks and selector validation via the \texttt{RoleRegistry} and per-function access
control modifiers. The CLI-side \texttt{CustodyGuard} catches mistakes early---wrong
contract address, wrong function selector, excessive value transfers---before wasting
gas on a transaction that would revert on-chain. Both layers must agree for a
transaction to succeed. Neither layer alone is sufficient: the on-chain layer cannot
prevent gas waste, and the CLI layer cannot prevent a malicious binary from bypassing it.
\end{criticalbox}

\subsection{Custody Policy per Job}

Each job definition includes a \texttt{CustodyPolicy} specifying the exact contracts,
selectors, value caps, and rate limits allowed.

\begin{table}[H]
\centering
\small
\begin{tabularx}{\textwidth}{l l l c}
\toprule
\textbf{Job} & \textbf{Allowed Destinations} & \textbf{Allowed Selectors} & \textbf{Value Cap} \\
\midrule
\rowcolor{keeper!8}
Oracle Updater & OracleManager & \texttt{updatePriceFeeds} & 0 \\
\rowcolor{keeper!8}
Funding Keeper & ClearingHouse & \texttt{settleFundingWindow} & 0 \\
\rowcolor{keeper!8}
Liq.\ Flagger & ClearingHouse & \texttt{flagPosition} & 0 \\
\rowcolor{operator!8}
Liq.\ Executor & ClearingHouse & \texttt{liquidatePosition} & 0 \\
\rowcolor{operator!8}
TP/SL Agent & ClearingHouse & \texttt{placeOrder} (REDUCE\_ONLY) & 0 \\
\rowcolor{cooperative!8}
Market Maker & Relay (off-chain) & commit, reveal & 0 \\
\rowcolor{cooperative!8}
ABM Agent & Relay + BinManager & commit, reveal, \texttt{recenterBins} & 0 \\
\rowcolor{managed!8}
GLV Manager & Glv & \texttt{deposit}, \texttt{withdraw}, \texttt{harvestPerformanceFee} & 0 \\
\bottomrule
\end{tabularx}
\caption{Custody policies per job. All jobs have a native value cap of 0 (no ETH/HYPE transfers). Cooperative jobs route through the relay for commit/reveal and do not submit on-chain transactions directly except for ABM bin recentering.}
\end{table}

%==============================================================================
\section{Job Registry (Local + On-Chain)}
%==============================================================================

The job registry exists in two layers: a local registry embedded in the CLI for
offline job discovery and validation, and an on-chain registry contract for
registration, staking, heartbeats, and reward distribution.

\subsection{Local Registry}

The local \texttt{JOB\_REGISTRY} mirrors the pattern established by
\texttt{strategy\_registry.py} in the agent runtime. It provides job metadata
for CLI commands (\texttt{hl jobs list}, \texttt{hl jobs info}) and validation
logic (\texttt{hl jobs run}) without requiring a chain connection.

\begin{contractbox}
\textbf{JOB\_REGISTRY}
\begin{lstlisting}[language=Python]
JOB_REGISTRY: Dict[str, JobDefinition] = {
    "oracle_updater": JobDefinition(
        job_id="oracle_updater",
        name="Oracle Updater",
        category=JobCategory.KEEPER,
        trigger=TriggerType.NEW_BLOCK,
        custody=CustodyPolicy(
            destinations=["OracleManager"],
            selectors=["updatePriceFeeds"],
            value_cap=0,
            max_txs_per_block=1,
        ),
        min_stake=10,
        requires_tee=False,
    ),
    "funding_keeper": JobDefinition(
        job_id="funding_keeper",
        name="Funding Keeper",
        category=JobCategory.KEEPER,
        trigger=TriggerType.NEW_BLOCK,
        custody=CustodyPolicy(
            destinations=["ClearingHouse"],
            selectors=["settleFundingWindow"],
            value_cap=0,
            max_txs_per_block=1,
        ),
        min_stake=10,
        requires_tee=False,
    ),
    "liq_flagger": JobDefinition(
        job_id="liq_flagger",
        name="Liquidation Flagger",
        category=JobCategory.KEEPER,
        trigger=TriggerType.ORACLE_UPDATE,
        custody=CustodyPolicy(
            destinations=["ClearingHouse"],
            selectors=["flagPosition"],
            value_cap=0,
            max_txs_per_block=5,
        ),
        min_stake=10,
        requires_tee=False,
    ),
    "liq_executor": JobDefinition(
        job_id="liq_executor",
        name="Liquidation Executor",
        category=JobCategory.OPERATOR,
        trigger=TriggerType.POSITION_FLAGGED,
        custody=CustodyPolicy(
            destinations=["ClearingHouse"],
            selectors=["liquidatePosition"],
            value_cap=0,
            max_txs_per_block=3,
        ),
        min_stake=50,
        requires_tee=False,
    ),
    "tpsl_agent": JobDefinition(
        job_id="tpsl_agent",
        name="TP/SL Agent",
        category=JobCategory.OPERATOR,
        trigger=TriggerType.ORACLE_UPDATE,
        custody=CustodyPolicy(
            destinations=["ClearingHouse"],
            selectors=["placeOrder"],
            value_cap=0,
            max_txs_per_block=10,
        ),
        min_stake=25,
        requires_tee=False,
    ),
    "market_maker": JobDefinition(
        job_id="market_maker",
        name="Market Maker",
        category=JobCategory.COOPERATIVE,
        trigger=TriggerType.ROUND_START,
        custody=CustodyPolicy(
            destinations=["Relay"],
            selectors=["commit", "reveal"],
            value_cap=0,
            max_txs_per_block=1,
        ),
        min_stake=100,
        requires_tee=True,
    ),
    "abm_agent": JobDefinition(
        job_id="abm_agent",
        name="ABM Agent",
        category=JobCategory.COOPERATIVE,
        trigger=TriggerType.ORACLE_UPDATE,
        custody=CustodyPolicy(
            destinations=["Relay", "BinManager"],
            selectors=[
                "commit", "reveal", "recenterBins",
            ],
            value_cap=0,
            max_txs_per_block=2,
        ),
        min_stake=100,
        requires_tee=True,
    ),
    "glv_manager": JobDefinition(
        job_id="glv_manager",
        name="GLV Manager",
        category=JobCategory.MANAGED,
        trigger=TriggerType.PERIODIC_AND_EVENT,
        custody=CustodyPolicy(
            destinations=["Glv"],
            selectors=[
                "deposit", "withdraw",
                "harvestPerformanceFee",
            ],
            value_cap=0,
            max_txs_per_block=2,
        ),
        min_stake=50,
        requires_tee=False,
    ),
}
\end{lstlisting}
\end{contractbox}

\subsection{On-Chain Registry Client}

The on-chain \texttt{JobRegistry} contract manages agent registration, staking,
heartbeat tracking, reward accrual, and slashing. The CLI interacts with it through
a \texttt{JobRegistryClient} abstraction that supports both mock and live
implementations.

\begin{contractbox}
\textbf{JobRegistryClient ABC}
\begin{lstlisting}[language=Python]
class JobRegistryClient(ABC):
    """Interface to the on-chain JobRegistry contract."""

    @abstractmethod
    def register(
        self,
        job_id: str,
        stake: int,
        attestation: bytes,
    ) -> str:
        """Register for a job with stake and optional
        TEE attestation. Returns transaction hash."""
        ...

    @abstractmethod
    def heartbeat(self, job_id: str) -> None:
        """Submit a heartbeat proving liveness."""
        ...

    @abstractmethod
    def claim_reward(self, job_id: str) -> int:
        """Claim accumulated rewards.
        Returns amount claimed in wei."""
        ...

    @abstractmethod
    def deregister(self, job_id: str) -> None:
        """Deregister and initiate unstaking cooldown."""
        ...
\end{lstlisting}
\end{contractbox}

\begin{contractbox}
\textbf{MockJobRegistryClient}
\begin{lstlisting}[language=Python]
class MockJobRegistryClient(JobRegistryClient):
    """For development before JobRegistry contract is
    deployed. Stores state in-memory. Heartbeats are
    no-ops. Rewards accrue at a fixed rate per call.
    Used in --dry-run mode and local development."""

    def __init__(self):
        self._registered: Dict[str, int] = {}
        self._rewards: Dict[str, int] = {}

    def register(self, job_id, stake, attestation):
        self._registered[job_id] = stake
        self._rewards[job_id] = 0
        return "0x" + "0" * 64  # mock tx hash

    def heartbeat(self, job_id):
        self._rewards[job_id] = (
            self._rewards.get(job_id, 0) + 1
        )

    def claim_reward(self, job_id):
        amount = self._rewards.get(job_id, 0)
        self._rewards[job_id] = 0
        return amount

    def deregister(self, job_id):
        self._registered.pop(job_id, None)
\end{lstlisting}
\end{contractbox}

\begin{contractbox}
\textbf{Web3JobRegistryClient}
\begin{lstlisting}[language=Python]
class Web3JobRegistryClient(JobRegistryClient):
    """Interacts with deployed JobRegistry contract via
    web3.py. Reads contract address and ABI from the
    deployments registry. Handles transaction signing,
    gas estimation, and confirmation waiting."""

    def __init__(
        self,
        rpc_url: str,
        registry_address: str,
        signer: Account,
    ):
        self._w3 = Web3(Web3.HTTPProvider(rpc_url))
        self._contract = self._w3.eth.contract(
            address=registry_address,
            abi=self._load_abi(),
        )
        self._signer = signer

    def register(self, job_id, stake, attestation):
        tx = self._contract.functions.register(
            job_id, attestation
        ).build_transaction({
            "value": stake,
            "from": self._signer.address,
        })
        signed = self._signer.sign_transaction(tx)
        tx_hash = self._w3.eth.send_raw_transaction(
            signed.raw_transaction
        )
        return tx_hash.hex()

    def heartbeat(self, job_id):
        tx = self._contract.functions.heartbeat(
            job_id
        ).build_transaction({
            "from": self._signer.address,
        })
        signed = self._signer.sign_transaction(tx)
        self._w3.eth.send_raw_transaction(
            signed.raw_transaction
        )

    def claim_reward(self, job_id):
        tx = self._contract.functions.claimReward(
            job_id
        ).build_transaction({
            "from": self._signer.address,
        })
        signed = self._signer.sign_transaction(tx)
        receipt = self._w3.eth.send_raw_transaction(
            signed.raw_transaction
        )
        # Parse reward amount from event logs
        return self._parse_reward_claimed(receipt)

    def deregister(self, job_id):
        tx = self._contract.functions.deregister(
            job_id
        ).build_transaction({
            "from": self._signer.address,
        })
        signed = self._signer.sign_transaction(tx)
        self._w3.eth.send_raw_transaction(
            signed.raw_transaction
        )
\end{lstlisting}
\end{contractbox}

%==============================================================================
\section{Cooperative Engine Bridge}
%==============================================================================

The \texttt{CooperativeEngine} bridges the job system with the existing TEE-work
cooperative market-making infrastructure. Rather than reimplementing the relay
polling, commit--reveal, and clearing feedback loop, it wraps the \texttt{AgentClient}
from the tee-work-llm repository.

\begin{contractbox}
\textbf{CooperativeEngine}
\begin{lstlisting}[language=Python]
class CooperativeEngine(JobEngine):
    """Engine for COOPERATIVE jobs.

    Wraps tee-work-llm's AgentClient, adding job lifecycle
    management (heartbeats, reward tracking, graceful
    shutdown) on top of the existing relay-based cooperative
    execution model.
    """

    def start(self, job_def: JobDefinition, config: JobConfig):
        from agent.client import AgentClient

        # Load the strategy (e.g., AvellanedaStoikovMM)
        strategy = load_strategy(config.strategy)

        # Initialize the AgentClient with job config
        self._client = AgentClient(
            agent_id=config.agent_id,
            strategy=strategy,
            relay_url=config.relay_url,
        )

        # TEE attestation if required
        if job_def.requires_tee:
            self._attest()

        # Start heartbeat background task
        self._heartbeat_task = asyncio.create_task(
            self._heartbeat_loop()
        )

        # Enter the cooperative execution loop
        # (poll relay -> tick -> seal -> commit -> reveal)
        self._client.run()

    def stop(self):
        self._heartbeat_task.cancel()
        self._client.shutdown()

    def heartbeat(self):
        self._registry_client.heartbeat(
            self._job_def.job_id
        )

    def status(self) -> JobStatus:
        return JobStatus(
            job_id=self._job_def.job_id,
            engine="CooperativeEngine",
            running=self._client.is_running(),
            uptime=self._uptime(),
            last_heartbeat=self._last_heartbeat,
            rewards=self._accumulated_rewards,
        )

    async def _heartbeat_loop(self):
        """Send heartbeats at regular intervals."""
        while True:
            await asyncio.sleep(30)
            self.heartbeat()
            self._last_heartbeat = time.time()
\end{lstlisting}
\end{contractbox}

\begin{designbox}[Shared Models Between Repositories]
The \texttt{CooperativeEngine} imports \texttt{AgentClient} from the tee-work-llm
repository. This import is safe because both repositories share identical data
models: \texttt{MarketSnapshot}, \texttt{StrategyDecision}, and \texttt{BaseStrategy}
are defined in tee-work-llm and consumed by the agent-cli.

Currently, these models are duplicated or imported via path manipulation. Long-term,
they should be extracted into a standalone \texttt{nunchi-sdk} package that both
repositories depend on. This eliminates version drift and makes the dependency
explicit in each project's \texttt{pyproject.toml}.

Until \texttt{nunchi-sdk} is published, the agent-cli lists tee-work-llm as a
development dependency and imports directly. The CI pipeline verifies model
compatibility by running cross-repository integration tests.
\end{designbox}

%==============================================================================
\section{File Structure}
%==============================================================================

The job integration code lives under \texttt{agent-cli/cli/jobs/}. Each module
has a single responsibility, matching the architecture sections above.

\begin{lstlisting}[language={}]
agent-cli/cli/jobs/
  __init__.py                  # Package init, exports command group
  commands.py                  # Typer command group (list, info,
                               #   register, run, status, claim,
                               #   deregister)
  registry.py                  # JOB_REGISTRY dict + JobDefinition
                               #   dataclass + JobCategory enum
  config.py                    # JobConfig dataclass (parsed from YAML)
  engines.py                   # JobEngine ABC + KeeperEngine +
                               #   CooperativeEngine + ManagedEngine +
                               #   JobEngineFactory
  events.py                    # ChainEvent dataclass + EventSubscriber
                               #   ABC + ChainEventSubscriber +
                               #   LogPollingSubscriber
  custody.py                   # CustodyPolicy dataclass + CustodyGuard
  strategy_interfaces.py       # KeeperStrategy + CooperativeStrategy +
                               #   ManagedStrategy ABCs + context models
  status.py                    # JobStatus dataclass + heartbeat/reward
                               #   tracking utilities
\end{lstlisting}

Integration with the existing CLI requires a two-line change to \texttt{cli/main.py}:

\begin{lstlisting}[language=Python]
# cli/main.py
from cli.jobs.commands import jobs_app

app.add_typer(jobs_app, name="jobs", help="Manage perpetual agent jobs")
\end{lstlisting}

%==============================================================================
\section{Migration and Phasing}
%==============================================================================

The CLI job integration rolls out in four phases, aligned with the migration plan
from the Perpetual Agent Jobs Specification.

\begin{table}[H]
\centering
\small
\begin{tabularx}{\textwidth}{l l X}
\toprule
\textbf{Phase} & \textbf{Infrastructure} & \textbf{CLI Capability} \\
\midrule
\rowcolor{gray!10}
Phase 1 & V1 Contracts &
\texttt{hl run} for free-form trading strategies. TEE-work scripts
(\texttt{run\_agent.py}, \texttt{run\_relay.py}) for cooperative market making.
No \texttt{hl jobs} command group yet. Keeper services run as ad-hoc cron jobs. \\
\rowcolor{info!10}
Phase 2 & L1 Testnet &
\texttt{hl jobs run oracle\_updater} and \texttt{hl jobs run liq\_flagger}
available as permissionless keeper jobs using \texttt{LogPollingSubscriber}
against V1 contracts on HyperEVM. \texttt{hl jobs run market\_maker} wraps
the tee-work \texttt{AgentClient} via \texttt{CooperativeEngine}.
\texttt{MockJobRegistryClient} used for registration and rewards (no on-chain
contract yet). \\
\rowcolor{success!10}
Phase 3 & L1 Mainnet &
All 8 jobs available via \texttt{hl jobs run}. \texttt{ChainEventSubscriber}
replaces \texttt{LogPollingSubscriber} for real-time event delivery.
\texttt{Web3JobRegistryClient} connects to the deployed \texttt{JobRegistry}
contract for on-chain registration, staking, heartbeats, and reward claims.
Full custody enforcement active on both CLI and contract layers. \\
\rowcolor{cooperative!8}
Phase 4 & Full Integration &
Agent Passports integrated into the registration flow. Dynamic job discovery
from the on-chain registry---the CLI queries the contract for available jobs
rather than relying solely on the local \texttt{JOB\_REGISTRY}. Passport-based
reputation scoring influences job assignment priority and reward multipliers. \\
\bottomrule
\end{tabularx}
\caption{Phased rollout of CLI job integration, aligned with infrastructure milestones.}
\end{table}

\begin{successbox}
The CLI job integration is considered complete when:
\begin{itemize}[leftmargin=2em]
  \item All 8 jobs can be started, monitored, and stopped via \texttt{hl jobs run}.
  \item The \texttt{JobEngineFactory} correctly routes each job category to its engine.
  \item Custody enforcement rejects transactions outside the job's policy at the CLI layer.
  \item Heartbeats are submitted automatically and visible via \texttt{hl jobs status}.
  \item Rewards can be claimed via \texttt{hl jobs claim} against the live \texttt{JobRegistry} contract.
  \item A new job type can be added by writing a strategy implementation and a registry entry, with no changes to the engine or routing infrastructure.
\end{itemize}
\end{successbox}

\end{document}
